\section{Anatomy of the shift handling}
\label{sec:shifts}

SHA-256's shifts and rotations are the structure that most cleanly
separates the two systems, and they are also the most common source of
confusion, because two different ``shifts'' live on two different axes.
We separate them, then show how each system handles each, and pin down
which one actually causes \binius{}'s linear verifier. (It is not the one
usually blamed.)

\begin{definition}[Two shift axes]
\label{def:two-shifts}
The \emph{entry-wise word shift} acts \emph{within} a word: the
$\Sigma/\sigma$ rotations and the logical shifts of SHA-256
($\ROTR$, $\SHR$ by fixed amounts on the $\wbit$ bits of one cell). The
\emph{cross-entry (wiring) shift} acts \emph{across} words: which
committed word supplies an operand, i.e. the dataflow between trace
positions.
\end{definition}

\subsection{\binius{}: the shift indicator and the monster sweep}
\label{sec:shifts-binius}

\binius{}'s Shift reduction realises both axes in one protocol, and it is
worth reading the verifier's final check to see exactly where the cost
lands. After the BitAnd reduction the verifier holds evaluation claims on
the operand multilinears $A,B,C$ at a random point. The Shift reduction
discharges them in two sumchecks and a product check.

\paragraph{The entry-wise shift is a constant-size indicator.}
The intra-word shift is captured by a \emph{shift indicator}
$\mathrm{shift\_ind}(i,j,s)$, the multilinear extension over the
constant-size $\log_2\wbit$-bit subcube of the Boolean predicate ``output
bit $i$ reads source bit $j$ under shift amount $s$.'' Its semantics are
the integer relations
\[
  \text{srl: } i+s=j, \quad
  \text{sll: } j+s=i, \quad
  \text{sra: } \min(i+s,\wbit{-}1)=j, \quad
  \text{rotr: } (i+s)\bmod\wbit = j,
\]
(plus $32$-bit lane variants), and it is evaluated in $O(\wbit)$ by short
$\sigma/\varphi$ recurrences contracted against the subspace-Lagrange
weights of the bit challenge. There are eight variants, so the per-output
work is a fixed $\mathrm{SHIFT\_VARIANT\_COUNT}\times\wbit = 8\times 64 =
512$-slot reduction. \emph{This is the entry-wise shift, and it is
$O(1)$ in $N$.}

\paragraph{The cross-entry wiring is the $O(N)$ ``monster.''}
The expensive part is the multilinear that encodes \emph{which committed
word feeds which operand}. The verifier evaluates it by sweeping over all
$N$ constraints: for each constraint, for each operand term, it
accumulates $\eq(r_{x'},\text{constraint})\cdot\eq(r_y,\text{word index})$
into the $512$ shift-variant slots. The two sumchecks reduce the bit and
shift-amount axes (constant size) and the word-index axis ($\log$ of the
word count) and collapse the witness side to a \emph{single} evaluation;
the final check is
\[
  \mathrm{eval} \;\stackrel{?}{=}\; \mathrm{witness\_eval}\cdot\mathrm{monster\_eval},
\]
where $\mathrm{monster\_eval}$ is purely public --- the $O(\wbit)$ shift
indicators times the $O(N)$ wiring sweep. The wiring sweep is the
verifier's hot loop; measured, it is $97$--$99.7\%$ of verify time
(\Cref{sec:results}).

\begin{observation}[It is the wiring, not the shifts]
\label{obs:wiring-not-shifts}
\binius{}'s linear verifier is caused by the cross-entry wiring sweep, not
by the entry-wise shifts. The shift \emph{operation} is succinct: the
indicator is $O(\wbit)$, evaluated once. The sweep is $O(N)$ only because
the per-constraint source-word selection is arbitrary public data with no
closed form. A common intuition --- ``shifts are why the verifier is
linear'' --- is contradicted by \zinc{}, which handles the identical
SHA-256 rotations with a succinct verifier (\Cref{sec:shifts-zinc}).
\end{observation}

\subsection{\zinc{}: shifts as free functional read-offs}
\label{sec:shifts-zinc}

In \zinc{}'s cell ring the entry-wise word shift is multiplication by a
fixed monomial $X^c$, and the key point is what this does at the
commitment layer. A committed cell exposes, as $\Ftwo$-linear read-offs,
\emph{any} linear functional of its bits; in particular it exposes both
projections $\psia$ (evaluate the bit-polynomial at $\alpha$) and $\psiz$
(the subspace-Lagrange weights of the Hadamard discharge) from one bound
object $a'$. A shift composes with these: $\psia(\SHR^c\, v)$ is a fixed
($X^c$-twisted) closed-form function of $\psia(v)$ up to a boundary
correction, folded into the $O(\log N)$ multipoint-evaluation step. So
the entry-wise shift costs the verifier nothing beyond reading a second
functional off the object it already opened, and the cross-entry dataflow
is the AIR's fixed next-row shift, checked in $O(\log N)$. There is no
per-constraint sweep to perform.

\subsection{The representational layer underneath}
\label{sec:shifts-repr}

\Cref{obs:wiring-not-shifts} says uniformity, not the shift itself, sets
verifier complexity. A second, finer point concerns \emph{how cheaply}
the entry-wise shift is realised, and it is a property of the word
representation rather than of uniformity:
\begin{itemize}
  \item \emph{$\Ftwo[X]$ cell ring} (\zinc{}): a shift is one scalar
  (monomial) multiply on the bound evaluation --- the cheapest form.
  \item \emph{bit-decomposed columns} (prime-field STARKs): a rotation is
  a free re-indexing of bit columns, at $\wbit$-fold trace width.
  \item \emph{packed flat-bit multilinear} (\binius{}): a shift is a
  hypercube-index transform handled by a constant-size shift reduction
  --- $O(1)$ per term, but not a single scalar multiply.
\end{itemize}
The asymptotics are the same in all three (a uniform shift touches only
the constant intra-word subcube), so a uniform AIR over any of these
representations has a succinct verifier. What differs is the
\emph{prover-side} cost of the entry-wise shift, and that is where the
$\Ftwo[X]$ representation has a structural advantage we quantify next:
in \binius{} the entry-wise shift is paid for by a sizeable Shift
reduction \emph{prover}, whereas in \zinc{} it is a free read-off.
